% GENERATED FILE. Edit canonical lessons, then run npm run generate:books.
% canonical-source-hash: fnv1a64:a3a81aabfafce9bf
% canonical-lessons: BN-C434-prokriti, BN-C434-poribesh, BN-C434-dushon, BN-C434-noshto-kora, BN-C434-songrokkhon-kora

\chapter{Nature and the Environment}
\label{ch:bn-nature-and-the-environment}
\hlchaptermodality{434}

\begin{chapteropening}
\textbf{By the end of this chapter:} \emph{I can talk about nature, pollution and what we must protect.}

Complete the last lesson of chapter 434: I can talk about nature, pollution and what we must protect.
\end{chapteropening}

\section[\texorpdfstring{\bn{প্রকৃতি} (prôkriti)}{প্রকৃতি (prôkriti)}]{\bn{প্রকৃতি} (prôkriti) — nature; the natural world}

\label{lesson:BN-C434-prokriti}



\begin{quote}
Before the new one: say the Bengali for to sway, then the Bengali for to babble (of water).
\end{quote}

\subsection*{You'll want to know: \bn{প্রকৃতি}}
\textbf{\bn{প্রকৃতি}} — \emph{prôkriti} — \textquotedblleft{}nature; the natural world\textquotedblright{}.

\begin{quote}
\bn{গ্রামে} \bn{এসে} \bn{মীরা} \bn{প্রকৃতির} \bn{শান্তি} \bn{পেল।}
\end{quote}

\begin{quote}
\emph{grāme ese mīrā prôkritir shānti pelô.}
\end{quote}

\begin{quote}
Coming to the village, Mira found the peace of nature.
\end{quote}

Also a person's nature, temperament.

\begin{cousinweb}
\textbf{\bn{প্রকৃতি}} (\emph{prôkriti}) is Sanskrit \emph{prakṛti}, \textquotedblleft{}the original or natural form of a thing, nature\textquotedblright{} (Monier-Williams), taken over whole.
\end{cousinweb}

\begin{grammarlens}[title={What you've built}]
One more word for getting around.
\end{grammarlens}

\subsection*{Your turn}
\begin{itemize}
\raggedright
  \item \emph{Say it:} \emph{prôkriti}
  \item \emph{Say it:} \emph{prôkriti}, once more
  \item \emph{Say it:} \emph{grāme ese mīrā prôkritir shānti pelô.} — Coming to the village, Mira found the peace of nature.
  \item \emph{Recall:} say the Bengali for to sway, then the Bengali for to babble (of water), then say \emph{prôkriti} again
\end{itemize}

\begin{tcolorbox}[breakable,colback=teal!4,colframe=teal!35!black,title={Before you move on}]
What does \bn{প্রকৃতি} mean? (\textquotedblleft{}Nature; the natural world\textquotedblright{}.) Say it once more.
\end{tcolorbox}
\section[\texorpdfstring{\bn{পরিবেশ} (pôribesh)}{পরিবেশ (pôribesh)}]{\bn{পরিবেশ} (pôribesh) — the environment; surroundings}

\label{lesson:BN-C434-poribesh}



\begin{quote}
Before the new one: say the Bengali for to babble (of water), then the Bengali for nature.
\end{quote}

\subsection*{You'll want to know: \bn{পরিবেশ}}
\textbf{\bn{পরিবেশ}} — \emph{pôribesh} — \textquotedblleft{}the environment; surroundings\textquotedblright{}.

\begin{quote}
\bn{কারখানার} \bn{ধোঁয়া} \bn{পরিবেশের} \bn{ক্ষতি} \bn{করে।}
\end{quote}

\begin{quote}
\emph{kārkhānār dhõyā pôribesher khôti kôre.}
\end{quote}

\begin{quote}
Factory smoke harms the environment.
\end{quote}

Also the surroundings of a place: \bn{শান্ত} \bn{পরিবেশ}, a calm setting.

\begin{cousinweb}
\textbf{\bn{পরিবেশ}} (\emph{pôribesh}) is Sanskrit \emph{pariveṣa}, \textquotedblleft{}a surrounding, a circle, a halo\textquotedblright{} (Monier-Williams); the sense \textquotedblleft{}the environment\textquotedblright{} is a modern use.
\end{cousinweb}

\begin{grammarlens}[title={What you've built}]
One more word for getting around.
\end{grammarlens}

\subsection*{Your turn}
\begin{itemize}
\raggedright
  \item \emph{Say it:} \emph{pôribesh}
  \item \emph{Say it:} \emph{pôribesh}, once more
  \item \emph{Say it:} \emph{kārkhānār dhõyā pôribesher khôti kôre.} — Factory smoke harms the environment.
  \item \emph{Recall:} say the Bengali for to babble (of water), then the Bengali for nature, then say \emph{pôribesh} again
\end{itemize}

\begin{tcolorbox}[breakable,colback=teal!4,colframe=teal!35!black,title={Before you move on}]
What does \bn{পরিবেশ} mean? (\textquotedblleft{}The environment; surroundings\textquotedblright{}.) Say it once more.
\end{tcolorbox}
\section[\texorpdfstring{\bn{দূষণ} (dūshôn)}{দূষণ (dūshôn)}]{\bn{দূষণ} (dūshôn) — pollution}

\label{lesson:BN-C434-dushon}



\begin{quote}
Before the new one: say the Bengali for nature, then the Bengali for the environment.
\end{quote}

\subsection*{You'll want to know: \bn{দূষণ}}
\textbf{\bn{দূষণ}} — \emph{dūshôn} — \textquotedblleft{}pollution\textquotedblright{}.

\begin{quote}
\bn{শহরে} \bn{গাড়ির} \bn{ধোঁয়ায়} \bn{দূষণ} \bn{হয়।}
\end{quote}

\begin{quote}
\emph{shôhôre gāṛir dhõyāy dūshôn hôy.}
\end{quote}

\begin{quote}
In the city car fumes cause pollution.
\end{quote}

\begin{cousinweb}
\textbf{\bn{দূষণ}} (\emph{dūshôn}) is Sanskrit \emph{dūṣaṇa}, \textquotedblleft{}spoiling, corrupting, vitiating\textquotedblright{} (Monier-Williams); the sense \textquotedblleft{}pollution\textquotedblright{} is a modern use.
\end{cousinweb}

\begin{grammarlens}[title={What you've built}]
One more word for getting around.
\end{grammarlens}

\subsection*{Your turn}
\begin{itemize}
\raggedright
  \item \emph{Say it:} \emph{dūshôn}
  \item \emph{Say it:} \emph{dūshôn}, once more
  \item \emph{Say it:} \emph{shôhôre gāṛir dhõyāy dūshôn hôy.} — In the city car fumes cause pollution.
  \item \emph{Recall:} say the Bengali for nature, then the Bengali for the environment, then say \emph{dūshôn} again
\end{itemize}

\begin{tcolorbox}[breakable,colback=teal!4,colframe=teal!35!black,title={Before you move on}]
What does \bn{দূষণ} mean? (\textquotedblleft{}Pollution\textquotedblright{}.) Say it once more.
\end{tcolorbox}
\section[\texorpdfstring{\bn{নষ্ট} \bn{করা} (nôshṭô kôrā)}{নষ্ট করা (nôshṭô kôrā)}]{\bn{নষ্ট} \bn{করা} (nôshṭô kôrā) — to spoil; to waste; to ruin}

\label{lesson:BN-C434-noshto-kora}



\begin{quote}
Before the new one: say the Bengali for the environment, then the Bengali for pollution.
\end{quote}

\subsection*{You'll want to know: \bn{নষ্ট} \bn{করা}}
\textbf{\bn{নষ্ট} \bn{করা}} — \emph{nôshṭô kôrā} — \textquotedblleft{}to spoil; to waste; to ruin\textquotedblright{}.

\begin{quote}
\bn{বন্যা} \bn{খেতের} \bn{ধান} \bn{নষ্ট} \bn{করল।}
\end{quote}

\begin{quote}
\emph{bônyā kheter dhān nôshṭô kôrlô.}
\end{quote}

\begin{quote}
The flood ruined the rice in the fields.
\end{quote}

Past (he/she …): \bn{সেই} \bn{বন্যা} \bn{ধান} \bn{নষ্ট} \bn{করল}, \emph{sei bônyā dhān nôshṭô kôrlô}.

\bn{নষ্ট} alone is spoiled, gone bad: \bn{খাবারটা} \bn{নষ্ট}, the food has gone off. \bn{সময়} \bn{নষ্ট} \bn{করা}, to waste time.

\begin{cousinweb}
\textbf{\bn{নষ্ট} \bn{করা}} (\emph{nôshṭô kôrā}) joins Sanskrit \emph{naṣṭa}, \textquotedblleft{}lost, perished, destroyed\textquotedblright{} (Monier-Williams), to \bn{করা}, to do.
\end{cousinweb}

\begin{grammarlens}[title={What you've built}]
One more word for getting around.
\end{grammarlens}

\subsection*{Your turn}
\begin{itemize}
\raggedright
  \item \emph{Say it:} \emph{nôshṭô kôrā}
  \item \emph{Say it:} \emph{nôshṭô kôrā}, once more
  \item \emph{Say it:} \emph{bônyā kheter dhān nôshṭô kôrlô.} — The flood ruined the rice in the fields.
  \item \emph{Recall:} say the Bengali for the environment, then the Bengali for pollution, then say \emph{nôshṭô kôrā} again
\end{itemize}

\begin{tcolorbox}[breakable,colback=teal!4,colframe=teal!35!black,title={Before you move on}]
What does \bn{নষ্ট} \bn{করা} mean? (\textquotedblleft{}To spoil; to waste; to ruin\textquotedblright{}.) Say it once more.
\end{tcolorbox}
\section[\texorpdfstring{\bn{সংরক্ষণ} \bn{করা} (sôngrôkkhôn kôrā)}{সংরক্ষণ করা (sôngrôkkhôn kôrā)}]{\bn{সংরক্ষণ} \bn{করা} (sôngrôkkhôn kôrā) — to conserve; to preserve}

\label{lesson:BN-C434-songrokkhon-kora}



\begin{quote}
Before the new one: say the Bengali for pollution, then the Bengali for to spoil.
\end{quote}

\subsection*{You'll want to know: \bn{সংরক্ষণ} \bn{করা}}
\textbf{\bn{সংরক্ষণ} \bn{করা}} — \emph{sôngrôkkhôn kôrā} — \textquotedblleft{}to conserve; to preserve\textquotedblright{}.

\begin{quote}
\bn{আমাদের} \bn{বন} \bn{আর} \bn{নদী} \bn{সংরক্ষণ} \bn{করতে} \bn{হবে।}
\end{quote}

\begin{quote}
\emph{āmāder bôn ār nôdī sôngrôkkhôn kôrte hôbe.}
\end{quote}

\begin{quote}
We must conserve our forests and rivers.
\end{quote}

Past (he/she …): \bn{তিনি} \bn{পুরনো} \bn{বই} \bn{সংরক্ষণ} \bn{করলেন}, \emph{tini purôno bôi sôngrôkkhôn kôrlen}.

\begin{cousinweb}
\textbf{\bn{সংরক্ষণ} \bn{করা}} (\emph{sôngrôkkhôn kôrā}) joins Sanskrit \emph{saṃrakṣaṇa}, \textquotedblleft{}guarding, keeping, preservation\textquotedblright{} (Monier-Williams), to \bn{করা}, to do.
\end{cousinweb}

\begin{grammarlens}[title={What you've built}]
One more word for getting around.
\end{grammarlens}

\subsection*{Your turn}
\begin{itemize}
\raggedright
  \item \emph{Say it:} \emph{sôngrôkkhôn kôrā}
  \item \emph{Say it:} \emph{sôngrôkkhôn kôrā}, once more
  \item \emph{Say it:} \emph{āmāder bôn ār nôdī sôngrôkkhôn kôrte hôbe.} — We must conserve our forests and rivers.
  \item \emph{Recall:} say the Bengali for pollution, then the Bengali for to spoil, then say \emph{sôngrôkkhôn kôrā} again
\end{itemize}

\begin{tcolorbox}[breakable,colback=teal!4,colframe=teal!35!black,title={Before you move on}]
What does \bn{সংরক্ষণ} \bn{করা} mean? (\textquotedblleft{}To conserve; to preserve\textquotedblright{}.) Say it once more.
\end{tcolorbox}
